% ================================================================
%  conteudo/material_metodos.tex
% ================================================================
\section{Material e Métodos}

\subsection{Material}

\begin{table}[H]
\caption{Materiais utilizados no experimento}
\centering
\begin{tabular}{cc}
\toprule
\textbf{Material} & \textbf{Especificação} \\
\midrule
Uvas                  & Variedade roxa \\
Água peptonada        & Concentração de 0,1\% \\
Ágar PDA              & Meio de cultura sólido \\
Ácido tartárico       & Ajuste de pH para 3,5 \\
Placas de Petri       & 6 unidades \\
Swabs                 & Estéreis \\
Béquer                & -- \\
Bastão de vidro       & -- \\
Bico de Bunsen        & -- \\
Estufa bacteriológica & -- \\
Autoclave             & Esterilização do meio \\
Gral e pistilo        & Maceração da amostra \\
\bottomrule
\end{tabular}
\end{table}

\subsection{Métodos}

O meio de cultura utilizado para o isolamento de micro-organismos foi o Ágar PDA (\textit{Potato Dextrose Agar}), previamente preparado no Laboratório de Biotecnologia da Escola de Química e Alimentos (EQA) da FURG. O preparo do meio foi realizado conforme as instruções de diluição do fabricante, utilizando água peptonada a 0,1\%, com ajuste de acidez para pH 3,5 por meio da adição de ácido tartárico. Posteriormente, o meio foi esterilizado em autoclave, a 121~°C e 1,1~atm, durante 15 minutos, sendo então vertido em placas de Petri estéreis.

A amostra empregada no experimento consistiu em um cacho de uvas adquirido no comércio local de Rio Grande -- RS. Para o isolamento dos microrganismos, cada integrante do grupo utilizou uma porção distinta da amostra, contemplando materiais provenientes da baga macerada, da casca e do engaço.

O experimento foi conduzido sob condições de assepsia, com higienização prévia das mãos e da bancada com álcool 70\%, mantendo-se o campo estéril com o auxílio da chama de um bico de Bunsen. Para a coleta dos microrganismos, foram utilizados swabs estéreis umedecidos em água peptonada a 0,1\%, a fim de favorecer a transferência e a viabilidade celular durante a inoculação.

Após a coleta, o material foi inoculado em placas contendo Ágar PDA por meio da técnica de esgotamento por estrias, com o objetivo de promover a separação física das células e possibilitar a obtenção de colônias isoladas. Em seguida, as placas devidamente identificadas foram incubadas em estufa bacteriológica a 30~°C por seis dias, permitindo o desenvolvimento das colônias para posterior observação macroscópica.